\documentclass[]{aiaa-LSP}% insert '[draft]' option to show overfull boxes
\usepackage{graphicx}
\usepackage{epstopdf}
\usepackage[spanish]{babel}
\usepackage[utf8]{inputenc}
\graphicspath{{Figuras/}}
	\addto\captionsspanish{\def\tablename{Tabla}}
\usepackage {hyperref}

 \title{Plantilla para Laboratorio de Sistemas de Propulsión}

 \author{
  Primer Autor%
    \thanks{Matrícula: 1234567.}
  \\ Segundo Autor%
   \thanks{Matrícula: 1234567.}\\
  {\normalsize\itshape
  Día:Lunes Hora:N1  Equipo:01}
 }

 % Data used by 'handcarry' option if invoked
 \AIAApapernumber{YEAR-NUMBER}
 \AIAAconference{Conference Name, Date, and Location}
 \AIAAcopyright{\AIAAcopyrightD{YEAR}}

 % Define commands to assure consistent treatment throughout document
 \newcommand{\eqnref}[1]{(\ref{#1})}
 \newcommand{\class}[1]{\texttt{#1}}
 \newcommand{\package}[1]{\texttt{#1}}
 \newcommand{\file}[1]{\texttt{#1}}
 \newcommand{\BibTeX}{\textsc{Bib}\TeX}

\begin{document}

\maketitle

\begin{abstract}
AQUÍ SE ESCRIBE EL RESUMEN
\end{abstract}

\section*{Nomenclatura}

\begin{tabbing}
  XXX \= \kill% this line sets tab stop
  $\rho$ \> Densidad, $kg/m^3$ \\
  $F$ \> Fuerza, N \\
 \end{tabbing}

\section{Introduccción}

Aquí se escribe la introducción.

\subsection{Subsección}

Algo de información 

\section{Ecuaciones}

Para insertar una ecuación y se muestra en la Ec.~(\ref{e:function}).

\begin{equation}
 \label{e:function}
D= \frac{1}{2}\rho V^2 A C_D
\end{equation}

La ecuación de la fuerza al arrastre es muy importante.\cite{Ali16}

\section{Gráficas}

Para insertar una figura, puede ser:
\begin{enumerate}
\item gráfica
\item imagen 
\end{enumerate}

En la fig. \ref{f:mag} se muestra un ejemplo de cómo se visualiza una imagen.

\begin{figure}[htb]% order of placement preference: here, top, bottom
\begin{center}
 \includegraphics[width=0.6\textwidth]{grafica}
 \caption{Aquí se escribe el texto.}
 \label{f:mag}
\end{center}
\end{figure}


\section{Tablas}

Para insertar una tabla se puede generar en la página www.tablesgenerator.com . Un ejemplo de cómo se visualiza una tabla, se muestra en la tabla \ref{t:ejemplo}


\begin{table}[h]
\centering
\caption{Nombre de la tabla.}
\label{t:ejemplo}
\begin{tabular}{cc}
Valor 1 & Valor 2 \\
1       & 2      
\end{tabular}
\end{table}

\section{Conclusiones}

Aquí se escriben sus conclusiones


\section*{Apéndice}

Aquí se escribe algún apéndice en caso de ser necesario.

\section*{Agradecimientos}
\c
Aquí se escribe en caso de querer agradecer. \cite{Napo12}

% Inicia la bibliografía
\bibliography{MiBiblio}
\bibliographystyle{aiaa}

\end{document}

% - Release $Name:  $ -
